\begin{abstract}
This study proposes to investigate the relationship between occupational exposure to artificial intelligence and labor market transitions in Brazil using the Continuous National Household Sample Survey. The research considers occupational changes, employment exits, and transitions into informality, using the public release of ChatGPT as a temporal reference for comparing periods. The design uses pairs of consecutive interviews, occupational exposure measures, and worker characteristics. The procedures include preparing inputs, mapping occupational classifications, constructing the panel, and estimating conditional associations while accounting for sampling weights, controls, and fixed effects. The theoretical framework distinguishes potential exposure, actual technology adoption, and observed employment outcomes. The analysis will examine the magnitude of the associations and their limitations, including changes in data collection and sample composition. This version is a provisional dissertation outline prepared for exercise ia-6.1 and presents no new empirical results. Its intended contribution is to organize a reproducible investigation of the topic without automatically attributing causality to observed differences. Conclusions will be developed after the author verifies the results and assesses the methodological limitations.
\end{abstract}